\section{Bayesian Receiver Design}
\label{sec:algorithm}

We derive the Gaussian updates and the local activity mixture under
\eqref{eq:working-posterior}, and then specify the finite-layer receiver.
The construction uses Gaussian projection and message division
\cite{minka2001}, the activity mixture in \cite[Sec.~III-B]{cai2026},
and receiver training motivated by \cite{wei2026}.
Learned calibration and residual feedback are distinguished from exact
local posterior moments. Let $J=KN$, $C=MN$, and assume
$\mathbf A_{\rm c}\ne\mathbf0$, $\sigma_{\rm c}^2>0$.

\subsection{Linear Estimation}

At layer $\ell$, the incoming Gaussian message has mean
$\mathbf M_\ell\in\mathbb C^{J\times C}$ and covariance $w_\ell\I_J$
per column, with $w_\ell>0$. Completing the square gives
\begin{align}
 \mathbf C_\ell
 &=\left(w_\ell^{-1}\I_J+
 \frac{\mathbf A_{\rm c}^{\mathsf H}\mathbf A_{\rm c}}{\sigma_{\rm c}^2}
 \right)^{-1},
 \label{eq:twc-linear-covariance}\\
 \mathbf M_{{\rm L},\ell}&=\mathbf M_\ell+\bm\Delta_\ell,
 \label{eq:twc-linear-mean}\\
 \bm\Delta_\ell&=w_\ell\mathbf A_{\rm c}^{\mathsf H}
 (\sigma_{\rm c}^2\I_Q+w_\ell\mathbf A_{\rm c}\mathbf A_{\rm c}^{\mathsf H})^{-1}
 (\mathbf Y-\mathbf A_{\rm c}\mathbf M_\ell).
 \label{eq:twc-linear-correction}
\end{align}
These expressions do not require row orthogonality. Replacing
$\mathbf C_\ell$ by $t_\ell\I_J$, where $t_\ell=\tr(\mathbf C_\ell)/J$,
is an isotropic Gaussian moment projection. Dividing the projected
Gaussian by the incoming message yields
\begin{equation}
 v_\ell^{(0)}=(t_\ell^{-1}-w_\ell^{-1})^{-1},\qquad
 \mathbf R_\ell^{(0)}=v_\ell^{(0)}
 \left(\frac{\mathbf M_{{\rm L},\ell}}{t_\ell}
       -\frac{\mathbf M_\ell}{w_\ell}\right).
 \label{eq:twc-linear-extrinsic}
\end{equation}
The eigenvalues of $\mathbf C_\ell$ lie in $(0,w_\ell]$, with at least
one strictly below $w_\ell$, so $0<t_\ell<w_\ell$.

Let $s_i$, $i=1,\ldots,r_A$, be the nonzero singular values of
$\mathbf A_{\rm c}$, where $r_A=\operatorname{rank}(\mathbf A_{\rm c})$.
Define
\begin{equation}
 f_\ell=\frac{J-r_A+\sum_{i=1}^{r_A}
 \sigma_{\rm c}^2/(\sigma_{\rm c}^2+w_\ell s_i^2)}{J},\qquad
 d_\ell=\frac1J\sum_{i=1}^{r_A}
 \frac{w_\ell s_i^2}{\sigma_{\rm c}^2+w_\ell s_i^2}.
 \label{eq:twc-linear-fractions}
\end{equation}
Then $f_\ell=t_\ell/w_\ell$ and $d_\ell=1-f_\ell$. Computing $d_\ell$
by its own sum avoids subtracting nearly equal quantities. Thus
\begin{equation}
 v_\ell^{(0)}=w_\ell f_\ell/d_\ell,\qquad
 \mathbf R_\ell^{(0)}=\mathbf M_\ell+\bm\Delta_\ell/d_\ell.
 \label{eq:twc-query-uncalibrated}
\end{equation}
For the unfolded receiver, introduce $\gamma_\ell=e^{c_{\gamma,\ell}}$
and use
\begin{equation}
 \xi_\ell=\gamma_\ell f_\ell/d_\ell,\qquad
 v_\ell=w_\ell\xi_\ell,\qquad
 \mathbf R_\ell=\mathbf M_\ell+(1+\xi_\ell)\bm\Delta_\ell.
 \label{eq:twc-query-calibrated}
\end{equation}
At $\gamma_\ell=1$, this reduces to \eqref{eq:twc-query-uncalibrated}.
Otherwise it modifies both the query mean and variance. This differs
from multiplying the posterior variance before message division in
\cite[eqs.~(8)--(10)]{wei2026}, which gives the ratio
$\gamma_\ell f_\ell/(1-\gamma_\ell f_\ell)$.

Even the uncalibrated query need not have independent Gaussian error.
For supplied physical parameters, set
$\mathbf L_\ell=w_\ell\mathbf A^{\mathsf H}
(\sigma^2\I_Q+w_\ell\mathbf A\mathbf A^{\mathsf H})^{-1}/d_\ell$.
Direct substitution gives
\begin{equation}
 \mathbf R_\ell^{(0)}-\mathbf X
 =(\mathbf L_\ell\mathbf A-\I_J)(\mathbf X-\mathbf M_\ell)
 +\mathbf L_\ell\mathbf W.
 \label{eq:twc-query-error}
\end{equation}
This motivates adapting the denoiser to receiver queries rather than
identifying every query with its AWGN pretraining model.

\subsection{Channel and Activity Denoising}

Reshape a user block of $\mathbf R_\ell$ as $r\in\mathbb C^D$,
$D=MN^2$. For the local statistical interpretation, use
\begin{equation}
 r=\alpha h+e,\qquad e\sim\CN(0,v\I_D),\qquad
 \alpha\sim\operatorname{Bernoulli}(\rho),\quad h\sim q_{\rm tr},
 \label{eq:twc-local-model}
\end{equation}
where $e$, $\alpha$, and $h$ are independent in this local model.
The hypothesis densities are
\begin{equation}
 q_0(r;v)=\CN(r;0,v\I_D),\qquad
 q_1(r;v)=\int\CN(r;h,v\I_D)q_{\rm tr}(dh).
 \label{eq:twc-local-densities}
\end{equation}
Gaussian convolution defines a smooth density even if $q_{\rm tr}$
is supported on a lower-dimensional set. Bayes' rule gives
\begin{equation}
 p^\star=\frac{\rho q_1(r;v)}{(1-\rho)q_0(r;v)+\rho q_1(r;v)},\quad
 \mu^\star=\E[h\mid r,\alpha=1],\quad
 \mathbf C_H^\star=\operatorname{Cov}(h\mid r,\alpha=1).
 \label{eq:twc-local-targets}
\end{equation}
For $x=\alpha h$, total expectation and covariance give
\begin{align}
 \E[x\mid r]&=p^\star\mu^\star,
 \label{eq:twc-mixture-mean}\\
 \operatorname{Cov}(x\mid r)&=p^\star\mathbf C_H^\star+
 p^\star(1-p^\star)\mu^\star\mu^{\star\mathsf H}.
 \label{eq:twc-mixture-covariance}
\end{align}
These are the mixture relations underlying
\cite[eqs.~(35)--(38)]{cai2026}.

Stack the real and imaginary parts of each user query as
$u_{k\ell}\in\mathbb R^{2D}$. The implemented layer outputs are
\begin{align}
 \mu_{k\ell}&=u_{k\ell}+a_\ell\frac{v_\ell}{2}
 s_{\theta_\ell}(u_{k\ell},v_\ell),\qquad a_\ell=e^{c_{\alpha,\ell}},
 \label{eq:twc-conditional-mean}\\
 b_{k\ell}&=\logit(\rho)+
 g_{\phi_\ell}(u_{k\ell},v_\ell,\mu_{k\ell},\rho),\qquad
 p_{k\ell}=\sigmoid(b_{k\ell}),
 \label{eq:twc-activity}\\
 \widehat x_{k\ell}&=p_{k\ell}\mu_{k\ell}.
 \label{eq:twc-estimate}
\end{align}
The factor $v_\ell/2$ is the variance per real coordinate. The scalar
$a_\ell$ calibrates denoising and is distinct from activity $\alpha_k$.
Score and classifier parameters are independent across layers and
shared across users within each layer. The classifier uses query and
conditional-estimate maps, their energies, and the public activity prior.
It is supervised as a posterior classifier on receiver queries; its
output is not assumed to be an exact AWGN likelihood ratio.
Reassemble \eqref{eq:twc-estimate} as $\widehat{\mathbf X}_\ell$
using \eqref{eq:mmv-indexing}.

\subsection{Extrinsic Message Update}

We form a feedback scale without a second-order score network using
the residual construction of \cite[eq.~(12)]{wei2026}, applied to the
effective channel:
\begin{equation}
 e_\ell=\frac{\|\mathbf Y-\mathbf A_{\rm c}\widehat{\mathbf X}_\ell\|_F^2
                   -QC\sigma_{\rm c}^2}
                  {C\|\mathbf A_{\rm c}\|_F^2}.
 \label{eq:twc-residual-scale}
\end{equation}
This is a feedback parameter, not the covariance in
\eqref{eq:twc-mixture-covariance}. For supplied physical parameters,
with $\mathbf E_\ell=\mathbf X-\widehat{\mathbf X}_\ell$, its numerator is
\begin{equation}
 \|\mathbf A\mathbf E_\ell\|_F^2+
 \|\mathbf W\|_F^2-QC\sigma^2+
 2\operatorname{Re}\langle\mathbf A\mathbf E_\ell,\mathbf W\rangle_F,
 \label{eq:twc-residual-expansion}
\end{equation}
where $\langle U,V\rangle_F=\tr(U^{\mathsf H}V)$.
Since the estimate depends on $\mathbf W$ through $\mathbf Y$,
the cross term need not have zero expectation and $e_\ell$ may be negative.

Let $\operatorname{softplus}(x)=\log(1+e^x)$ and define
\begin{align}
 T_\ell&=\eta v_\ell,\qquad
 z_\ell=e^{c_{\beta,\ell}}
 [T_\ell\operatorname{softplus}(e_\ell/T_\ell)+\epsilon_s],
 \label{eq:twc-positive-scale}\\
 q_\ell&=\epsilon_r+(q_{\max}-\epsilon_r)
 \tanh\!\left(\frac{z_\ell}{(q_{\max}-\epsilon_r)v_\ell}\right),
 \qquad \tau_\ell=v_\ell q_\ell.
 \label{eq:twc-feedback-ratio}
\end{align}
Here $\eta,\epsilon_s>0$ and $0<\epsilon_r<q_{\max}<1$ are fixed.
Consequently $0<q_\ell<1$. Dividing a Gaussian with mean
$\widehat{\mathbf X}_\ell$ and assigned variance $\tau_\ell$ by the
incoming Gaussian with mean $\mathbf R_\ell$ and variance $v_\ell$ gives
\begin{equation}
 \widetilde{\mathbf M}_{\ell+1}
 =\frac{\widehat{\mathbf X}_\ell-q_\ell\mathbf R_\ell}{1-q_\ell},
 \qquad \widetilde w_{\ell+1}=\frac{v_\ell q_\ell}{1-q_\ell}.
 \label{eq:twc-feedback-proposal}
\end{equation}
The restriction $\tau_\ell<v_\ell$ ensures positive outgoing precision;
it is not a general inequality for non-Gaussian posterior variances.

We then damp the Gaussian messages in natural parameters \cite{vehtari2020}:
\begin{align}
 \delta_\ell&=\min\{1,w_{\rm ref}/\widetilde w_{\ell+1}\},
 \label{eq:twc-damping-weight}\\
 w_{\ell+1}^{-1}&=(1-\delta_\ell)w_\ell^{-1}
                   +\delta_\ell\widetilde w_{\ell+1}^{-1},
 \label{eq:twc-damping-precision}\\
 \mathbf M_{\ell+1}&=w_{\ell+1}
 \left[(1-\delta_\ell)\frac{\mathbf M_\ell}{w_\ell}
       +\delta_\ell\frac{\widetilde{\mathbf M}_{\ell+1}}
                          {\widetilde w_{\ell+1}}\right].
 \label{eq:twc-damping-mean}
\end{align}
The reference variance $w_{\rm ref}$ is fixed at initialization.
Positive precision is preserved. The variance-dependent weight in
\eqref{eq:twc-damping-weight} is our implementation choice and does not
provide a convergence guarantee.
